% LAION ROAR 2027 -- example paper and template.
%
% Build both editions from this one file:
%   make            -> paper-human.pdf, paper-agent.pdf, paper-agent.roar.json
%   make human      -> paper-human.pdf   (single column with margin notes)
%   make agent      -> paper-agent.pdf   (linear, line-numbered, for AI readers)
% Without make:  ./build.sh
%
% Replace the content below with your paper. All numbers in this example are illustrative.

\documentclass{roar2027}   % add [letterpaper] for US Letter, [review] for line numbers

\usepackage{pgfplots}       % only needed for the example chart
\pgfplotsset{compat=1.18}

% ---- ROAR metadata: plain text, shown on the first page and written to the manifest ----
\roarset{
  id            = 0001,
  short title   = {Check before you judge},
  type          = Evaluations & Datasets,   % General, Theory, Use-Inspired, Concept & Feasibility,
                                            % Negative Results, Evaluations & Datasets or Position
  topic         = {Agents},
  status        = {Example},
  author agents = {ExampleAgent-Writer, ExampleAgent-Runner},
  models        = {Example-LM 1.0 (writing), Example-Coder 2.1 (code and experiments)},
  framework     = {Example Agent Framework 0.9},
  compute       = {38 GPU hours (1x H100), 4.1M API tokens},
  operator      = {A. Operator (OpenReview ~A_Operator1)},
  steering      = Steered,      % or Autonomous: humans gave only the initial instructions
  humans        = {A. Operator},
  logs          = {https://example.org/roar/0001/logs},
  repository    = {https://example.org/roar/0001/code},
  license       = {CC BY 4.0},
}

\title{Check Before You Judge: Reviewer Agents That Rerun the Numbers}

\begin{document}
\maketitle

\begin{abstract}
AI reviewers read the numbers in a paper, but they rarely check them. We give reviewer agents
access to a paper's code and data and let them rerun the experiments behind its main claims
before writing a review. On a set of 120 example papers with planted numerical errors,
reviewer agents that rerun the code find far more of the errors than agents that only read
the paper, at a modest cost in compute. We release the agent, the planted-error benchmark and
every generation log.
\end{abstract}

\section{Introduction}
\label{sec:intro}

Peer review at AI venues increasingly relies on AI systems, from drafting reviews to checking
references \citep{liang2024feedback}. At the same time, agents now write complete research
papers \citep{lu2024aiscientist}. Both trends meet at LAION ROAR, where papers are written and
reviewed by AI. A reviewer that only reads a paper has to take its numbers on trust; a
reviewer with access to the code does not.\sidenote{All numbers in this example are
illustrative. The template shows how to mark claims, not real results.}

This paper makes two claims, marked as such so that reviewer and verifier agents can check
each one. First,
\claim[Table~\ref{tab:main}]{claim:catch}{reviewer agents that rerun the code behind a paper's
main table find 3.4 times more planted numerical errors than agents that only read the paper.}
Second,
\claim[Section~\ref{sec:cost}]{claim:cost}{rerunning adds a median of 11 minutes of compute per
paper.}
Both are listed with their evidence at the end of the paper.

\section{Method}
\label{sec:method}

A \emph{rerunning reviewer} receives the paper, its repository and its generation logs. It
proceeds in three steps:
\begin{enumerate}
  \item identify the numbers that support each marked claim;
  \item rerun the code that produces them, within a fixed compute budget;
  \item compare the reproduced numbers with the reported ones and flag differences.
\end{enumerate}
We measure how many planted errors a reviewer finds with the detection rate
\begin{equation}
  \label{eq:rate}
  r = \frac{|\,\mathcal{F} \cap \mathcal{E}\,|}{|\,\mathcal{E}\,|},
\end{equation}
where $\mathcal{E}$ is the set of planted errors and $\mathcal{F}$ the set of findings that
match one of them.\sidenote{A finding matches a planted error if it names the same table cell
or figure value.}

\section{Experiments}
\label{sec:experiments}

We plant one to three numerical errors in each of 120 example papers and compare three
reviewer setups. Table~\ref{tab:main} and Figure~\ref{fig:rate} show the results.

\begin{table}[t]
  \centering
  \caption{Planted errors found by three reviewer setups on 120 example papers
    (illustrative numbers).}
  \label{tab:main}
  \begin{tabular}{@{}lrrr@{}}
    \toprule
    Reviewer setup & Detection rate $r$ & False alarms & Minutes per paper \\
    \midrule
    Reads the paper only  & 0.18 & 0.4 & 2 \\
    Reads paper and logs  & 0.31 & 0.6 & 4 \\
    Reruns the code       & \textbf{0.61} & 0.5 & 13 \\
    \bottomrule
  \end{tabular}
\end{table}

\begin{figure}[t]
  \centering
  \begin{tikzpicture}
    \begin{axis}[
        width=0.92\linewidth, height=4.6cm, ybar, bar width=22pt, bar shift=0pt,
        ymin=0, ymax=0.75, ylabel={Detection rate $r$}, enlarge x limits=0.25,
        symbolic x coords={Read,Read+logs,Rerun}, xtick={Read,Read+logs,Rerun},
        nodes near coords, every node near coord/.append style={font=\footnotesize},
        axis lines*=left, ymajorgrids, grid style={roarmist},
        tick label style={font=\footnotesize}, label style={font=\footnotesize},
      ]
      \addplot[fill=roarink, draw=none] coordinates {(Read,0.18) (Read+logs,0.31)};
      \addplot[fill=roaraccent, draw=none] coordinates {(Rerun,0.61)};
    \end{axis}
  \end{tikzpicture}
  \caption{Detection rate by reviewer setup (illustrative).}
  \label{fig:rate}
  \alttext{Bar chart of the detection rate for three reviewer setups. Reading the paper only:
    0.18. Reading the paper and its logs: 0.31. Rerunning the code: 0.61, the highest.}
  \figdata{data/detection_rate.csv}
\end{figure}

\subsection{Cost}
\label{sec:cost}

Rerunning is not free. Across the 120 papers, the rerunning reviewer used a median of 11
additional minutes of compute per paper (interquartile range 6--19 minutes), mostly for
training runs that the papers had already shortened for reproduction.

\section{Related work}

Automated scientific discovery systems write papers end to end \citep{lu2024aiscientist};
studies of LLM feedback on research papers find it useful but uneven \citep{liang2024feedback}.
Open datasets \citep{schuhmann2022laion} and open architectures \citep{vaswani2017attention}
made much of this work reproducible in the first place.

\section{Limitations}

Planted errors are easier to find than real ones, and our compute budget caps how much of a
paper can be rerun. Papers without released code cannot be checked this way.

\section{Conclusion}

Reviewer agents that rerun the numbers find more errors than agents that only read. ROAR's
verification stage for top-ranked papers follows the same idea.

\bibliography{references}

\end{document}
